\paragraph{3.1. Low-degree polynomials of Haar random unitaries.}
As discussed in the proof overview, we require bounds on the operator norm (as well as normalized trace) of low-degree polynomials of random unitaries and these follow from known results in free probability theory. Here we explain these connections and also prove some auxillary lemmas needed for the proof of Theorem 4 and Theorem 3.
Let $z_{\mathbf i}$ denote the non-commutative monomial $z_{i_1}z_{i_2}\cdots z_{i_d}$ for a $d$-tuple $\mathbf i=(i_1,\ldots,i_d)\in[t]^d$ and let $p(z_1,\ldots,z_t)$ be a non-commutative polynomial in the variables $z_1,\ldots,z_t$. We are interested in computing the operator norm $\|\cdot\|_{\mathrm{op}}$ and the normalized trace $\operatorname{tr}_N$ of the polynomial $p(z_1,\ldots,z_t)$ (or its higher moments) when substituting $N\times N$ Haar random unitaries for the variables $z_i$.
As explained previously, the theory of free probability gives us tools that allow us to compute the above in the limit $N\to\infty$. In particular, Voiculescu [35] showed that the (normalized) trace of polynomials in Haar random unitaries and their conjugates converges to the trace of the same polynomial evaluated on certain infinite-dimensional operators called Haar unitaries that satisfy a non-commutative notion of independence called free independence. This was strengthened by Collins and Male [16] who showed that such convergence also holds for the operator norm. A short primer on free probability is given in Appendix A.1, but for now one can think of $\mathcal A$ as a self-adjoint algebra of bounded linear operators on a Hilbert space and $\varphi$ as a trace functional for such operators in the statement given below.
\paragraph{Theorem 8 ([35,16]).}
Let $p(z_1,\ldots,z_{2t})$ be a non-commutative polynomial in $\mathbb R\langle z_1,\ldots,z_{2t}\rangle$. If $U_1,\ldots,U_t$ are $N\times N$ Haar random unitaries, then almost surely,
\[\lim_{N\to\infty}\operatorname{tr}_N[p(U_1,\ldots,U_t,U_1^*,\ldots,U_t^*)]=\varphi[p(u_1,\ldots,u_t,u_1^*,\ldots,u_t^*)],\]
\[\lim_{N\to\infty}\|p(U_1,\ldots,U_t,U_1^*,\ldots,U_t^*)\|_{\mathrm{op}}=\|p(u_1,\ldots,u_t,u_1^*,\ldots,u_t^*)\|,\]
where $u_1,\ldots,u_t$ are free Haar unitaries in a $C^*$-probability space $(\mathcal A,\varphi)$ and $\|\cdot\|$ is the norm for the underlying $C^*$-algebra.
Using the above result it suffices to consider free Haar unitaries in a $C^*$-probability space to compute the operator norm and trace of polynomials of random unitaries. For a non-commutative polynomial $p(z_1,\ldots,z_t)=\sum_{|\mathbf i|\leq d}c_{\mathbf i}z_{\mathbf i}$, denoting by $\|p\|_2=(\sum_{|\mathbf i|\leq d}|c_{\mathbf i}|^2)^{1/2}$, one can show the following easily using techniques from free probability.
\paragraph{Lemma 9.}
Let $p(z_1,\ldots,z_t)=\sum_{|\mathbf i|\leq d}c_{\mathbf i}z_{\mathbf i}$ be a non-commutative degree-$d$ polynomial in $\mathbb R\langle z_1,\ldots,z_t\rangle$ and $u_1,\ldots,u_t$ be free Haar unitaries in a $C^*$-probability space $(\mathcal A,\varphi)$. Then,
\[\varphi[p(u_1,\ldots,u_t)(p(u_1,\ldots,u_t))^*]=\|p\|_2^2.\]
\begin{proof}[Proof of Lemma 9 (Appendix A.2)]
Writing $u_{\mathbf i}^*=(u_{\mathbf i})^*$ for a tuple $\mathbf i$ and using linearity of $\varphi$, we have that
\[\varphi[p(u_1,\ldots,u_t)(p(u_1,\ldots,u_t))^*]=\sum_{|\mathbf i|,|\mathbf j|\leq d}c_{\mathbf i}c_{\mathbf j}\varphi(u_{\mathbf i}u_{\mathbf j}^*).\]
From Proposition 13, the term $\varphi(u_{\mathbf i}u_{\mathbf j}^*)$ is 1 iff $u_{\mathbf i}u_{\mathbf j}^*$ reduces to identity in the free group $F_t$ with generators $u_1,\ldots,u_t$. For the right-hand side above, this only happens when $\mathbf i=\mathbf j$ and thus these are the only non-zero terms. Thus,
\[\varphi[p(u_1,\ldots,u_t)(p(u_1,\ldots,u_t))^*]=\sum_{|\mathbf i|\leq d}|c_{\mathbf i}|^2.\]
\end{proof}
The above implies that $\operatorname{tr}_N[p(U_1,\ldots,U_t)(p(U_1,\ldots,U_t))^*]$ converges to $\|p\|_2^2$ almost surely as $N\to\infty$.
In particular, we will rely on the following result that follows from the work of Kemp and Speicher [23] who consider the operator norm of homogeneous polynomials evaluated on free $R$-diagonal operators, a class that includes free Haar unitaries as well. We also remark that a bound where the right-hand side below is worse by a multiplicative $O(d^{1/2})$ factor also follows from the work of Haagerup [22] who proved it in another context, predating even the introduction of free probability theory. Haagerup considered the more general case of polynomials in both $u_i$'s and $u_i^*$'s.
\paragraph{Theorem 10 ([23]).}
Let $p(z_1,\ldots,z_t)=\sum_{|\mathbf i|=d}c_{\mathbf i}z_{\mathbf i}$ be a homogeneous non-commutative degree-$d$ polynomial in $\mathbb R\langle z_1,\ldots,z_t\rangle$ and $u_1,\ldots,u_t$ be free Haar unitaries in a $C^*$-probability space. Then,
\[\|p(u_1,\ldots,u_t)\|\leq\sqrt{e(d+1)}\cdot\|p\|_2,\]
where the left-hand side denotes the norm in the underlying $C^*$-algebra.
We shall need to extend the above bound to non-homogeneous polynomials. Let $p(z_1,\ldots,z_t)=\sum_{|\mathbf i|\leq d}c_{\mathbf i}z_{\mathbf i}$ and let $p_k(z_1,\ldots,z_t)=\sum_{|\mathbf i|=k}c_{\mathbf i}z_{\mathbf i}$ denote the degree-$k$ homogeneous part of $p$. Writing $p_k=p_k(u_1,\ldots,u_t)$ for $0\leq k\leq d$ and $p=p(u_1,\ldots,u_t)$, it follows from the triangle inequality, Theorem 10, and Cauchy--Schwarz, that
\begin{align*}
\|p\|&\leq\sum_{k=0}^d\|p_k\|\leq\sum_{k=0}^d\sqrt{e(k+1)}\|p_k\|_2\\
&\leq\sqrt e\left(\sum_{k=0}^d(k+1)\right)^{1/2}\left(\sum_{k=0}^d\|p_k\|_2^2\right)^{1/2}\leq\sqrt e(d+1)\cdot\|p\|_2.
\end{align*}
Thus, we essentially get the same bound as in the homogeneous case, at the expense of an additional $O(d^{1/2})$ factor.
Collecting all the above we have the following as a direct consequence:
\paragraph{Theorem 11.}
Let $p(z_1,\ldots,z_t)=\sum_{|\mathbf i|\leq d}c_{\mathbf i}z_{\mathbf i}$ be a non-commutative degree-$d$ polynomial in $\mathbb R\langle z_1,\ldots,z_t\rangle$ and $U_1,\ldots,U_t$ be independent $N\times N$ Haar random unitaries. Then, the following holds almost surely,
\[\lim_{N\to\infty}\operatorname{tr}_N[p(U_1,\ldots,U_t)(p(U_1,\ldots,U_t))^*]=\|p\|_2^2,\]
and
\[\lim_{N\to\infty}\|p(U_1,\ldots,U_t)\|_{\mathrm{op}}\leq\sqrt e(d+1)\cdot\|p\|_2.\]
Moreover, the factor $(d+1)$ in the operator norm bound can be improved to $\sqrt{d+1}$ if additionally $p$ is also assumed to be homogeneous.
Based on the above theorem, we prove the following key lemma which captures the polar decomposition strategy mentioned in the earlier proof overview (Section 1.2). This will serve as the key ingredient in the proof of Theorem 3 and Theorem 4.
\paragraph{Lemma 12.}
Let $p$ be a non-commutative degree-$d$ polynomial in $\mathbb R\langle y_1,\ldots,y_m,z_1,\ldots,z_t\rangle$ given by
\[p(y_1,\ldots,y_m,z_1,\ldots,z_t)=\sum_{i=1}^my_iq_i(z_1,\ldots,z_t)+q_0(z_1,\ldots,z_t).\]
Then, for every $\delta>0$, there exist an integer $N$ and $N\times N$ unitaries $V_1,\ldots,V_m,W_1,\ldots,W_t$ such that
\[\|p(V_1,\ldots,V_m,W_1,\ldots,W_t)\|_{\mathrm{op}}\geq\frac1{\sqrt e(d+1)}\sum_{i=1}^m\|q_i\|_2-\delta.\]
Moreover, the factor in front can be improved to $(e(d+1))^{-1/2}$ if $p$ is homogeneous as well.
We note that the definition of completely bounded norm allows us to plug in matrices of arbitrarily large dimensions, so the exact dependence on $\delta$ is not relevant for our purposes since we can always make it arbitrarily small.
\begin{proof}
For an arbitrary integer $N$, let us pick independent $N\times N$ Haar random unitaries $W_1,\ldots,W_t$ which we substitute for the variables $z_1,\ldots,z_t$, respectively, and let $M_i=q_i(W_1,\ldots,W_t)$ be the corresponding random matrices. Then, for any tuple of matrices $V_1,\ldots,V_m$ that we could substitute for the variables $y_1,\ldots,y_m$, we have that
\[p(V_1,\ldots,V_m,W_1,\ldots,W_t)=\sum_{i=1}^mV_iM_i+M_0.\]
Theorem 11 and union bound imply that as $N\to\infty$, with probability 1 all the following events simultaneously hold:
\begin{itemize}
\item $\|M_i\|_{\mathrm{op}}\leq\sqrt e(d+1)\cdot\|q_i\|_2$ for each $i$,
\item $\operatorname{tr}_N(M_i^*M_i)=\|q_i\|_2^2$ for each $i$, where $\operatorname{tr}_N(M)$ is the normalized trace.
\end{itemize}
To show that the operator norm must be large, let us fix a sufficiently large $N$ and a choice of $N\times N$ unitaries $W_1,\ldots,W_t$ such that $M_i$ satisfies $\|M_i\|_{\mathrm{op}}\leq\sqrt e(d+1)\cdot\|q_i\|_2+\epsilon$ and $\operatorname{tr}_N(M_i^*M_i)\geq\|q_i\|_2^2-\epsilon$ for each $0\leq i\leq m$, where $\epsilon$ can be made arbitrarily small by increasing $N$. For $0\leq i\leq m$, let $M_i=U_iP_i$ be the left polar decomposition of $M_i$, where $U_i$ is a unitary matrix and $P_i$ is a positive semidefinite matrix.
We select the tuple of unitary matrices $V_1,\ldots,V_m$ that we substitute for the variables $y_1,\ldots,y_m$ to be $V_i=U_0U_i^*$ for $i\in[m]$. With this we have that $\|p(V_1,\ldots,V_m,W_1,\ldots,W_t)\|_{\mathrm{op}}$ is at least
\begin{align*}
\left\|M_0+\sum_{i=1}^mV_iM_i\right\|_{\mathrm{op}}
&=\left\|U_0P_0+\sum_{i=1}^mU_0U_i^*U_iP_i\right\|_{\mathrm{op}}\\
&=\left\|U_0P_0+\sum_{i=1}^mU_0P_i\right\|_{\mathrm{op}}\\
&=\left\|P_0+\sum_{i=1}^mP_i\right\|_{\mathrm{op}}\\
&\geq\operatorname{tr}_N\left(P_0+\sum_{i=1}^mP_i\right)\geq\operatorname{tr}_N\left(\sum_{i=1}^mP_i\right),
\end{align*}
where the last equality follows since the operator norm is unitarily invariant and the last two inequalities follow from the positive semidefiniteness of the $P_i$'s.
For every positive semidefinite matrix $P$, we have that $\operatorname{tr}_N(P)\geq\operatorname{tr}_N(P^2)/\|P\|_{\mathrm{op}}$. Hence,
\[\|p(V_1,\ldots,V_m,W_1,\ldots,W_t)\|_{\mathrm{op}}\geq\sum_{i=1}^m\frac{\operatorname{tr}_N(P_i^2)}{\|P_i\|_{\mathrm{op}}}.\]
By our choice of $M_i$, we have that $\operatorname{tr}_N(P_i^2)=\operatorname{tr}_N(M_i^*M_i)\geq\|q_i\|_2^2-\epsilon$ and $\|P_i\|_{\mathrm{op}}=\|M_i\|_{\mathrm{op}}\leq\sqrt e(d+1)\|q_i\|_2+\epsilon$. Since $\epsilon$ can be made arbitrarily small by increasing $N$, it follows that
\[\|p(V_1,\ldots,V_m,W_1,\ldots,W_t)\|_{\mathrm{op}}\geq\frac1{\sqrt e(d+1)}\sum_{i=1}^m\|q_i\|_2-\delta,\]
for large enough $N$. The improved bound for the homogeneous case follows directly by plugging the bound of Theorem 11 into the above proof.
\end{proof}
